\chapter{Conclusioni e sviluppi futuri}

\section{Sintesi e limiti dello studio}
Il presente elaborato ha analizzato l'impatto di due forme di degradazione controllata del dataset di training sulle prestazioni di classificazione di diverse architetture convoluzionali.
I dataset scelti sono Oxford-IIIT Pet e CIFAR-100, mentre le due tipologie di degradazione sono il macro-drop, che riduce uniformemente gli esempi di ciascuna classe, e il micro-drop, che rimuove del tutto alcune classi.

I risultati hanno mostrato con costanza come il macro-drop non produca una degradazione significativa delle prestazioni, 
mentre il micro-drop comprometta progressivamente il Macro F1-score al crescere della percentuale di classi rimosse, con un impatto relativo simile tra i due dataset ma livelli assoluti più critici su CIFAR-100, 
verosimilmente a causa della maggiore complessità intrinseca del task. In entrambi i casi, 
EfficientNet si è confermata l'architettura con le prestazioni assolute più elevate, mentre ResNet18, addestrata soltanto su Oxford-IIIT Pet, ha mostrato le maggiori criticità sia in termini di prestazioni assolute sia di stabilità tra le run.

A partire da questi risultati, è stato adottato un approccio a reti prototipiche per recuperare la capacità di riconoscere le classi escluse dal training tramite un numero ridotto di esempi di supporto. 
I risultati ottenuti su Oxford-IIIT Pet, valutati su $30$ episodi indipendenti, hanno mostrato un recupero parziale ma apprezzabile delle classi non viste, a fronte di un peggioramento moderato delle prestazioni sulle classi note, 
più marcato per alcune architetture rispetto ad altre.

I risultati vanno letti alla luce di alcuni limiti dello studio. Su Oxford-IIIT Pet ciascuna configurazione è stata valutata su due soli seed, mentre su CIFAR-100 sono state addestrate due sole architetture, su un singolo seed, e la valutazione prototipica è stata condotta su un singolo episodio. Le percentuali di micro-drop non sono direttamente confrontabili tra i due dataset, poiché su Oxford-IIIT Pet sono calcolate sulle sole razze canine. Il set di validazione, utilizzato per l'early stopping, include anche le classi rimosse dal training, e l'F1-score sulle classi note risente in parte dei campioni delle classi rimosse, assegnati come falsi positivi alle classi note. Infine, la struttura gerarchica è sfruttata soltanto attraverso il backbone condiviso: le prestazioni della head \emph{coarse} non sono state analizzate e non è stato effettuato un confronto con un classificatore \emph{flat}.

\section{Sviluppi futuri}
Diversi aspetti del presente lavoro si prestano a ulteriori approfondimenti. Un primo passo naturale consiste nell'estendere a CIFAR-100 la stessa profondità sperimentale adottata per Oxford-IIIT Pet, 
addestrando le architetture mancanti e ripetendo la valutazione del modello prototipico su un numero di episodi comparabile, in modo da rendere il confronto tra i due dataset pienamente omogeneo.

Un secondo sviluppo riguarda l'aumento del numero di seed indipendenti per ciascuna configurazione, così da poter stimare con maggiore accuratezza la variabilità dei risultati e distinguere in modo più netto gli effetti sistematici 
da quelli dovuti al rumore di inizializzazione.

Infine, un'ulteriore direzione di ricerca potrebbe riguardare la valutazione del modello prototipico secondo lo scenario di \emph{few-shot puro}, 
in cui la classificazione avviene esclusivamente tra le classi non viste, isolando così la capacità di apprendimento few-shot del modello dall'influenza delle classi già note; 
un confronto diretto tra questo scenario e quello \emph{Generalized} adottato in questo lavoro permetterebbe di quantificare più precisamente l'effetto, sulle classi non viste, 
della competizione con i prototipi delle classi note.

\section{Considerazioni finali}
Il presente elaborato ha affrontato un aspetto spesso trascurato nella valutazione dei modelli di classificazione: 
le loro prestazioni a fronte del degrado del dataset di training, distinguendo esplicitamente tra una riduzione uniforme degli esempi e una rimozione mirata di intere classi. 
I risultati ottenuti mostrano che queste due forme di degradazione hanno un impatto profondamente diverso sulle prestazioni, e che tale impatto non è uniforme tra le diverse architetture: 
EfficientNet ha ottenuto le prestazioni assolute più elevate in entrambi gli scenari testati, mentre un'architettura meno performante come ResNet18 ha mostrato una fragilità significativa sia in termini di prestazioni sia di stabilità.

L'adozione di un approccio a reti prototipiche ha inoltre dimostrato che è possibile recuperare, almeno parzialmente, 
la capacità di riconoscere le classi escluse dal training senza ricorrere a un nuovo addestramento completo del modello, sebbene con un costo non trascurabile in termini di accuratezza rispetto a quella ottenuta sulle classi apprese in modo tradizionale. 
Questo risultato conferma il potenziale del few-shot learning come strumento pratico per scenari in cui il dataset di training è incompleto o soggetto ad aggiornamenti frequenti, una condizione tutt'altro che rara 
in contesti reali.

Nel complesso, lo studio condotto offre indicazioni utili sia dal punto di vista della scelta architetturale in presenza di dataset potenzialmente incompleti, 
sia come base metodologica per l'integrazione di tecniche di few-shot learning in pipeline di classificazione già esistenti.